\scanpage{24}
\mspage{11}
c'est une fonction localement constante sur $\mathcal S$.

Fixons nous un type numérique $(g,\nu)$\inserted{\uline{anabélien} $(2g+\nu\geq3)$}, et considérons, pour $\mathcal S$ variable, le groupoïde fibré
\[
(\unsure{24})\qquad\mathcal S\longmapsto\widehat{\mathcal M}_{g,\nu}(\mathcal S)=\begin{gathered}\text{MD-courbes}\\\text{relatives sur }\mathcal S,\text{ de}\\\text{type numérique}\\\text{égal à }(g,\nu)\end{gathered}
\]
On a alors le vraiment mirifique théorème suivant :

\margin{On suppose $2g+\nu\geq3$ (cas anabélien)}
\uline{Théorème} de \uline{Mumford-Deligne} : Le

\figure{\begin{tikzpicture}
\node[anchor=north west,text width=0.9\linewidth,align=left,inner sep=4pt] (T) at (0,0) {
groupoïde fibré $\mathcal S\mapsto\widehat{\mathcal M}_{g,\nu}(\mathcal S)$ sur $(\mathrm{Sch})$ (plus gén\textsuperscript{l}, sur les multiplicités schématiques $\ldots$) est représentable par une multiplicité schématique $\widehat{\mathcal M}_{g,\nu}$, qui est \uline{lisse} et \uline{propre} sur $\operatorname{Spec}\mathbb Z$; \deleted{à fibres géom. connexes.}\inserted{D'autre part $\unsure{\mathcal M^!_{g,\nu}|\mathcal M'_{g,\nu}}$ est un}
};
\draw[line width=0.5pt] (T.north west)--(T.south west);
\end{tikzpicture}}
\mspage{23}
